%=====================================================================
\chapter{Container Orchestration with Kubernetes}\label{ch:kubernetes}
%=====================================================================
\locator{3}{Part III --- Containers and Lightweight Virtualization}

\begin{outcomes}
By the end of this chapter you will be able to:
\begin{tight}
  \item \textbf{Explain} the control loop and \textbf{distinguish} declarative
        from imperative management.
  \item \textbf{Name} the control-plane and node components and \textbf{trace}
        what each does when a Deployment is created.
  \item \textbf{Describe} a Pod, a Deployment and a Service, and \textbf{say}
        which problem each one solves.
  \item \textbf{Write} requests and limits for a workload, and \textbf{predict}
        the scheduling and eviction behaviour that follows.
  \item \textbf{Diagnose} a workload that will not schedule or will not stay
        running, from the resource arithmetic and the probe configuration.
\end{tight}
\end{outcomes}

\begin{prereq}
\chref{containers} (namespaces, cgroups, and above all that CPU limits throttle
while memory limits kill --- Kubernetes inherits both) and \chref{docker}
(images, digests, and the OCI runtime stack the kubelet drives).
\end{prereq}

\begin{keyterms}
orchestration $\bullet$ declarative $\bullet$ reconciliation $\bullet$ control
loop $\bullet$ desired state $\bullet$ API server $\bullet$ etcd $\bullet$
scheduler $\bullet$ controller manager $\bullet$ kubelet $\bullet$ kube-proxy
$\bullet$ Pod $\bullet$ sidecar $\bullet$ ReplicaSet $\bullet$ Deployment
$\bullet$ Service $\bullet$ ClusterIP $\bullet$ Ingress $\bullet$ request
$\bullet$ limit $\bullet$ QoS class $\bullet$ eviction $\bullet$ liveness probe
$\bullet$ readiness probe $\bullet$ taint $\bullet$ toleration
\end{keyterms}

\section{The Problem Containers Created}

\chref{docker} gave us an artefact that starts in a tenth of a second and costs
ten megabytes. \chref{containers} showed a host could hold two thousand of
them. Neither says anything about the questions that follow immediately:

Which host should this container run on? What happens when it dies at three in
the morning? When the host dies? How does traffic find it, given that its
address changes every restart? How do you replace all sixty copies with a new
version without dropping a request --- and undo that in forty seconds when it
is wrong?

Answering those by hand is possible for ten containers and impossible for a
thousand. \term{Orchestration} is the automation of those answers, and
Kubernetes is the one that won.

\begin{definitionbox}[Declarative Management and the Control Loop]
You do not tell Kubernetes what to \emph{do}. You tell it what you \emph{want},
and it works continuously to make reality match.

\smallskip
A \term{controller} is a loop that:
\begin{tight}
  \item reads the \term{desired state} --- what you declared;
  \item observes the \term{actual state} --- what exists;
  \item takes one step to reduce the difference;
  \item repeats, for ever.
\end{tight}

\smallskip
This is \term{reconciliation}, and it is the single idea the whole system is
built from. ``Run five copies'' is not a command executed once; it is a
condition continuously enforced. Kill one and a controller notices the
difference and creates another --- not because anything handled a failure, but
because four is not five.
\end{definitionbox}

\begin{conceptbox}[Why Declarative Beats Imperative Here]
An imperative system executes instructions: \emph{start a container},
\emph{stop that one}. It must therefore anticipate every way reality can
diverge --- a node rebooting, a network partition healing, a half-finished
rollout --- and have an instruction for each. The failure modes multiply
faster than anyone can enumerate them.

\smallskip
A declarative system has one failure mode: the difference between desired and
actual is not shrinking. Every problem reduces to that, and every recovery is
the same loop running again. It is why a Kubernetes cluster can be left alone
over a weekend and why \texttt{kubectl apply} is safe to run twice.

\smallskip
The cost is that you give up knowing exactly \emph{when} things happen. The
system converges; it does not promise a schedule. Operators used to imperative
tooling find this unnerving, and the answer is to watch the difference, not the
steps.
\end{conceptbox}

\section{What the Cluster Is Made Of}

\dsfig{%
\begin{tikzpicture}[font=\scriptsize,
  bx/.style={rounded corners=2pt, line width=0.8pt, align=center,
             inner sep=3pt, minimum width=2.6cm, minimum height=0.62cm,
             font=\scriptsize},
  cp/.style={bx, draw=primary, fill=primary!10, text=primary},
  nd/.style={bx, draw=greenhl, fill=greenhl!8, text=greenhl!45!black},
  st/.style={bx, draw=goldhl, fill=goldhl!12, text=goldhl!45!black},
  pod/.style={bx, draw=secondary, fill=secondary!8, text=secondary!80!black,
              minimum width=1.6cm, minimum height=0.5cm},
  ar/.style={-{Stealth[length=1.8mm]}, line width=0.8pt},
  note/.style={font=\scriptsize\itshape, text=neutral, align=left}
]
  % --- control plane ---
  \node[rounded corners=3pt, draw=primary!45, fill=primary!3,
        minimum width=6.0cm, minimum height=3.3cm] at (0,1.5) {};
  \node[font=\scriptsize\bfseries, text=primary] at (0,2.95) {Control plane};
  \node[cp] (api) at (0,2.3)  {\textbf{API server}};
  \node[st] (etcd) at (0,1.5) {etcd \ \ {\tiny the only state}};
  \node[cp] (sch) at (-1.5,0.6) {scheduler};
  \node[cp] (ctl) at (1.5,0.6)  {controller\\manager};
  \draw[ar, primary!60] (api) -- (etcd);
  \draw[ar, primary!60] (sch) -- (sch |- api.south);
  \draw[ar, primary!60] (ctl) -- (ctl |- api.south);

  % --- nodes ---
  \foreach \i/\x in {1/6.2, 2/10.4}{
    \node[rounded corners=3pt, draw=greenhl!45, fill=greenhl!3,
          minimum width=3.7cm, minimum height=3.3cm] at (\x,1.5) {};
    \node[font=\scriptsize\bfseries, text=greenhl!45!black] at (\x,2.95)
         {Node \i};
    \node[nd, minimum width=3.1cm] at (\x,2.3) {kubelet};
    \node[nd, minimum width=3.1cm] at (\x,0.35) {kube-proxy};
    \node[pod] at (\x-0.8,1.45) {Pod};
    \node[pod] at (\x+0.8,1.45) {Pod};
    \draw[ar, greenhl!60] (\x,2.0) -- (\x,1.75);}
  \draw[ar, primary!55] (api.east) to[out=0,in=180] (4.3,2.3);
  \draw[ar, primary!55] (3.1,2.1) to[out=0,in=180] (8.5,2.3);

  \node[note, anchor=west] at (-3.2,-0.9)
       {\textbf{API server} --- the only thing that writes to etcd. Everything
        else talks to it, and nothing talks to anything else.\\
        \textbf{etcd} --- the cluster's entire state, in one consistent
        key-value store. Lose it and you have lost the cluster.\\
        \textbf{scheduler} --- decides which node an unassigned Pod belongs on.
        It only writes that decision down; it does not start anything.\\
        \textbf{controller manager} --- the reconciliation loops: Deployments,
        ReplicaSets, nodes, endpoints.\\
        \textbf{kubelet} --- the only component that touches containers. Reads
        what should be on its node, drives containerd, reports back.\\
        \textbf{kube-proxy} --- programmes the node's routing so a Service
        address reaches a healthy Pod.};
\end{tikzpicture}}%
{A cluster. The arrows all point at the API server, which is the architectural
claim: components do not coordinate with each other, they each reconcile
against one shared record.}{k8s-components}

\begin{conceptbox}[Trace a Deployment Through the Loops]
\texttt{kubectl apply} a Deployment asking for three replicas. Nothing is
started by the thing that handled your request.

\begin{tight}
  \item The \textbf{API server} validates it and writes it to \textbf{etcd}.
        Your command returns. Nothing is running.
  \item The \textbf{Deployment controller} notices a Deployment with no
        ReplicaSet and creates one.
  \item The \textbf{ReplicaSet controller} notices a ReplicaSet wanting three
        Pods and zero existing, and creates three Pod objects --- with no node
        assigned.
  \item The \textbf{scheduler} notices three unassigned Pods, scores the nodes
        on free resources and constraints, and writes a node name into each.
        It still has not started anything.
  \item Each \textbf{kubelet} notices a Pod assigned to its node, pulls the
        image and asks containerd to run it.
  \item The \textbf{endpoints controller} notices Pods that match a Service's
        selector and adds their addresses; \textbf{kube-proxy} programmes the
        routing.
\end{tight}

\smallskip
Six independent loops, none of which called another. Each noticed a difference
and reduced it. Now kill a Pod: step 3's controller sees two where three are
wanted, and the same sequence runs from there. \emph{Failure recovery is not a
feature; it is the ordinary behaviour of the system.}
\end{conceptbox}

\section{Pods, Deployments and Services}

\rowcolors{2}{white}{lightbg}
\begin{longtable}{@{}L{2.6cm}L{5.4cm}L{6.7cm}@{}}
\toprule
\rowcolor{primary!15}
\textbf{Object} & \textbf{What it is} & \textbf{The problem it solves} \\
\midrule
\endhead
\term{Pod} & One or more containers sharing a network namespace, IP address
and storage volumes. The smallest schedulable unit & Some containers must be
co-located and share \texttt{localhost} --- an application and its log shipper,
or a \term{sidecar} proxy. You almost never create one directly \\
\term{ReplicaSet} & Keeps $n$ identical Pods running & The reconciliation loop
for ``how many''. You almost never create one directly either \\
\term{Deployment} & Manages ReplicaSets to roll out and roll back versions &
Changing the image without an outage. It creates a second ReplicaSet, shifts
Pods across gradually, and keeps the old one so a rollback is one command \\
\term{Service} & A stable virtual IP and DNS name in front of a changing set of
Pods & Pod addresses change on every restart. A \term{ClusterIP} does not,
and kube-proxy keeps its backend list current \\
\term{Ingress} & HTTP routing from outside the cluster to Services & A Service
is reachable inside the cluster; this is how the world gets in, with host and
path rules and TLS \\
\bottomrule
\end{longtable}

\begin{alertbox}[A Pod is not a container, and the difference bites]
A Pod is the unit of \emph{scheduling} and of \emph{IP address}. Every
container in it runs on the same node, shares one network namespace --- so they
reach each other on \texttt{localhost} and cannot both bind port 8080 --- and
can share volumes.

\smallskip
The error this produces is putting an application and its database in one Pod
because they belong together. They do not: they scale differently, fail
differently and are updated differently, and now they cannot be scheduled
separately or scaled separately. The test is whether the two must be on the
same machine and share a lifetime. A log shipper must. A database must not.
\end{alertbox}

\section{Requests, Limits and What the Scheduler Does With Them}

Every container may declare a \term{request} and a \term{limit} for CPU and
memory. They do completely different jobs and conflating them is the commonest
source of Kubernetes capacity trouble.

\begin{definitionbox}[Request and Limit]
The \term{request} is what the \textbf{scheduler} uses. It reserves that much
on the node, and a node whose remaining allocatable capacity is less than the
request will not be chosen. The request is a claim on \emph{placement} and is
reserved whether or not the container uses it.

\smallskip
The \term{limit} is what the \textbf{kernel} enforces, through the cgroups of
\chref{containers}. CPU above the limit is throttled; memory above it is an OOM
kill.

\smallskip
The gap between them is where a workload may burst, and where the cluster is
oversubscribed. A Pod requesting 100m of CPU with a 2-core limit is scheduled
as if it were tiny and may consume twenty times that --- from capacity the
scheduler has promised to somebody else.
\end{definitionbox}

\begin{worked}[Why This Pod Will Not Schedule]
A cluster has four nodes of 16 cores and 64 GB. A team reports that their new
Deployment sits in \texttt{Pending} and asks for more nodes. Decide whether
they need any.

\smallskip
\textbf{What is declared.} 20 replicas, each requesting 2 cores and 8 GB.

\smallskip
\textbf{Step 1 --- what the nodes actually offer.} Allocatable capacity is not
total capacity: the kubelet, the container runtime and the operating system are
reserved first. Typically about 10\% of CPU and 15\% of memory:
\[ \text{per node} \approx 16 \times 0.90 = 14.4 \text{ cores}, \qquad
   64 \times 0.85 = 54.4 \text{ GB} \]
\[ \text{cluster} = 4 \times 14.4 = 57.6 \text{ cores}, \qquad
   4 \times 54.4 = 217.6 \text{ GB} \]

\smallskip
\textbf{Step 2 --- what is already running.} Existing workloads request 30
cores and 120 GB. Remaining:
\[ 57.6 - 30 = 27.6 \text{ cores} \qquad 217.6 - 120 = 97.6 \text{ GB} \]

\smallskip
\textbf{Step 3 --- what the new Deployment asks for.}
\[ 20 \times 2 = 40 \text{ cores} \qquad 20 \times 8 = 160 \text{ GB} \]
Both exceed what is left. It cannot schedule, and the team is right that
something must change.

\smallskip
\textbf{Step 4 --- how many \emph{will} fit?} The scheduler places Pods one at
a time, so some will run. The binding constraint is memory:
\[ \left\lfloor \frac{97.6}{8} \right\rfloor = 12 \text{ Pods} \]
but Pods cannot be split across nodes, so compute it per node. If the existing
load is even, each node has $54.4 - 30 = 24.4$ GB free, giving
$\lfloor 24.4/8 \rfloor = 3$ Pods per node $= \mathbf{12}$. Eight stay
\texttt{Pending} --- which matches what they are seeing.

\smallskip
\textbf{Step 5 --- are the requests honest?} This is the question nobody asks,
and it is usually the answer. Measured over the previous week the same
application used 0.3 cores and 1.2 GB at the 95th percentile. The request of 2
cores and 8 GB is roughly \textbf{seven times} its actual demand.

\smallskip
Set requests to p95 plus a margin --- say 0.5 cores and 2 GB --- and limits
somewhat above:
\[ 20 \times 0.5 = 10 \text{ cores} \qquad 20 \times 2 = 40 \text{ GB} \]
Against 27.6 cores and 97.6 GB free, all twenty schedule with room to spare.

\smallskip
\textbf{Step 6 --- the recommendation.} No new nodes. The cluster was not full;
it was \emph{reserved}. A request is a promise the scheduler keeps whether the
workload uses it or not, so inflated requests buy hardware that then sits idle
--- the server sprawl of \chref{what}, re-created inside a cluster that was
bought to prevent it.
\end{worked}

\begin{conceptbox}[QoS Classes: Who Gets Evicted First]
From the requests and limits, Kubernetes derives a \term{QoS class} that
decides what happens when a node runs out of memory.

\begin{tight}
  \item \textbf{Guaranteed} --- every container has limits equal to requests.
        Evicted last.
  \item \textbf{Burstable} --- requests set, limits higher or absent. Evicted
        in the middle, worst offender against its request first.
  \item \textbf{BestEffort} --- nothing set at all. Evicted first, always.
\end{tight}

\smallskip
So omitting requests and limits is not neutral --- it is choosing to be killed
first. This surprises teams who leave them out to ``keep it simple'', and the
symptom arrives weeks later as a workload that dies only when the cluster is
busy.
\end{conceptbox}

\begin{conceptbox}[Liveness and Readiness Are Not the Same Question]
Two probes, routinely confused, with opposite consequences.

\smallskip
A \term{readiness} probe asks \emph{should traffic go here?} Failing it removes
the Pod from its Service's endpoints. Nothing is killed; the Pod is simply left
alone until it recovers. Use it for warm-up, for a lost database connection,
for shedding load.

\smallskip
A \term{liveness} probe asks \emph{is this process broken beyond recovery?}
Failing it \textbf{kills the container}. Use it only for genuine deadlock.

\smallskip
The classic outage: a liveness probe that hits an endpoint which also checks
the database. The database slows, every replica fails liveness, Kubernetes
kills all of them, they restart, reconnect in a thundering herd and fail again.
A readiness probe there would have removed them from rotation and let them
recover. \emph{If in doubt, use readiness.}
\end{conceptbox}

\begin{pitfall}
\begin{tight}
  \item \textbf{Setting requests from the limit, or from a guess.} The
        scheduler reserves the request. Inflated requests waste the cluster
        exactly as oversized virtual machines waste a host.
  \item \textbf{Omitting requests entirely.} That is BestEffort, which is
        evicted first. Not setting them is a decision.
  \item \textbf{Using a liveness probe where readiness is meant.} Liveness
        kills. The cascading restart above is the single most common
        self-inflicted Kubernetes outage.
  \item \textbf{Setting CPU limits aggressively.} CPU throttling is invisible
        in most dashboards and shows up as latency. Many teams set CPU requests
        and no CPU limit deliberately; check \texttt{throttled\_usec} before
        deciding.
  \item \textbf{Putting unrelated containers in one Pod.} They now share a
        lifetime, a node and a scaling decision.
  \item \textbf{Using \texttt{:latest} in a manifest.} The rollout is not
        reproducible and a rollback may not roll anything back
        (\chref{docker}).
  \item \textbf{Treating Kubernetes as a hypervisor.} It schedules containers
        on a shared kernel. Every isolation caveat in \chref{containers} still
        applies, and \chref{lightweight} is what to do about it.
\end{tight}
\end{pitfall}

\begin{lab}[A Cluster on One Laptop]
Builds a real cluster with \texttt{kind} --- a real API server, scheduler and
kubelet; only the node count is a pretence --- and produces each failure mode
above on purpose. Expect forty-five minutes.

\begin{lstlisting}[language=bash]
#!/usr/bin/env bash
# ch10_kind.sh -- reconciliation, scheduling and probes, observed.
set -euo pipefail

# --- 1. A three-node cluster in containers on this machine. ------------
kind create cluster --name vss --config - <<'YAML'
kind: Cluster
apiVersion: kind.x-k8s.io/v1alpha4
nodes:
  - role: control-plane
  - role: worker
  - role: worker
YAML
kubectl get nodes -o wide
# Each "node" is a container (Chapter 9) running a kubelet. Containers
# all the way down.

# --- 2. Declare a desired state and watch six loops reach it. ----------
kubectl create deployment web --image=nginx:1.27 --replicas=3
kubectl get deploy,rs,pods -o wide
kubectl get events --sort-by=.lastTimestamp | tail -12
# Read the events bottom-up: scheduled, pulling, created, started.

# --- 3. Reconciliation: delete a Pod and time the replacement. ---------
kubectl get pods -l app=web
kubectl delete pod -l app=web --field-selector status.phase=Running \
  --wait=false | head -1
sleep 3 ; kubectl get pods -l app=web
# Nothing "handled" the deletion. Two is not three, so a loop acted.

# --- 4. A Service: a stable address in front of changing Pods. --------
kubectl expose deployment web --port=80 --target-port=80
kubectl get svc web
kubectl get endpoints web          # the Pod IPs, maintained by a controller
kubectl delete pod -l app=web --wait=false | head -1
sleep 5 ; kubectl get endpoints web   # different Pod IPs, same Service IP

# --- 5. Make a Pod unschedulable, and read why. -----------------------
kubectl create deployment big --image=nginx:1.27
kubectl set resources deployment big --requests=cpu=8,memory=32Gi
kubectl get pods -l app=big
kubectl describe pod -l app=big | sed -n '/Events/,$p' | tail -6
# "0/3 nodes are available: Insufficient cpu" -- the worked example,
# reported by the scheduler itself.

# --- 6. QoS classes, derived from what you declared. ------------------
kubectl run guaranteed --image=nginx:1.27 \
  --overrides='{"spec":{"containers":[{"name":"c","image":"nginx:1.27",
  "resources":{"requests":{"cpu":"100m","memory":"64Mi"},
  "limits":{"cpu":"100m","memory":"64Mi"}}}]}}'
kubectl run besteffort --image=nginx:1.27
sleep 5
kubectl get pod guaranteed besteffort \
  -o custom-columns=NAME:.metadata.name,QOS:.status.qosClass

# --- 7. The cascading-restart outage, on purpose. ---------------------
# A liveness probe pointing at something that will fail.
kubectl create deployment fragile --image=nginx:1.27 --replicas=3
kubectl patch deployment fragile --type=json -p='[{"op":"add",
  "path":"/spec/template/spec/containers/0/livenessProbe",
  "value":{"httpGet":{"path":"/nope","port":80},"periodSeconds":5}}]'
sleep 45
kubectl get pods -l app=fragile      # RESTARTS climbing on every replica
kubectl describe pod -l app=fragile | grep -A2 'Liveness' | head -6
# Now imagine /nope were a real endpoint that checks a slow database.

# --- 8. Tear down. ----------------------------------------------------
kind delete cluster --name vss
\end{lstlisting}

\textbf{What to notice.} Step 3 is the whole chapter: nothing detected a
failure or ran a recovery procedure. A controller compared two numbers and
acted. Step 7 is the chapter's warning, reproduced in forty-five seconds.
\end{lab}

\begin{checkpoint}
\begin{tightnum}
  \item A cluster of six nodes has 12 allocatable cores and 48 GB each.
        Existing workloads request 40 cores and 190 GB. A Deployment of 15
        replicas requesting 1 core and 6 GB each is created. How many schedule,
        and what is the binding resource?
  \item A team sets no requests or limits anywhere ``to keep things simple''.
        Name the QoS class this produces and describe the symptom they will
        eventually see, including when.
  \item Explain why pointing a liveness probe at an endpoint that checks a
        shared database can take down every replica at once, and what to use
        instead.
\end{tightnum}
\end{checkpoint}

\begin{chaptersummary}
\begin{tight}
  \item Kubernetes is declarative: you state desired state and controllers
        reconcile towards it for ever. Failure recovery is not a feature but
        the ordinary behaviour of a loop comparing two numbers.
  \item Everything talks to the API server and nothing talks to anything else.
        etcd holds the entire cluster state.
  \item Creating a Deployment runs six independent loops; none calls another.
  \item A Pod is the unit of scheduling and of IP address. Put containers
        together only if they must share a node and a lifetime.
  \item Requests are for the scheduler and are reserved whether used or not;
        limits are enforced by cgroups, throttling CPU and killing on memory.
  \item QoS class follows from requests and limits, and decides eviction
        order. Declaring nothing means being evicted first.
  \item Readiness removes traffic; liveness kills the container. Using liveness
        where readiness was meant is the commonest self-inflicted outage.
\end{tight}
\end{chaptersummary}

\begin{reviewq}
  \item \textbf{(MCQ)} The component that decides which node a Pod runs on,
        and then does nothing further, is: (a)~the kubelet; (b)~the scheduler;
        (c)~kube-proxy; (d)~the API server.
  \item \textbf{(MCQ)} A container with no requests or limits is in QoS class:
        (a)~Guaranteed; (b)~Burstable; (c)~BestEffort; (d)~none.
  \item \textbf{(Short)} State the control loop in four steps and explain why
        it makes \texttt{kubectl apply} safe to run twice.
  \item \textbf{(Short)} Explain the difference between a request and a limit,
        naming which component enforces each.
  \item \textbf{(Short)} Give a workload for which two containers genuinely
        belong in one Pod, and one where they do not, with the reason.
  \item \textbf{(Applied)} A 10-node cluster of 32 cores and 128 GB each is
        reported ``full'' at 40\% measured utilisation. Explain how both facts
        can be true, describe what you would measure, and set out the change
        you would propose and its risk.
  \item \textbf{(Applied)} During a database slowdown, every replica of three
        different services restarted repeatedly and the outage lasted far
        longer than the database problem. Explain the mechanism, and write the
        probe configuration you would have had in place.
\end{reviewq}
